\vssub
\subsection{~气海过程}
\vsssub
\subsubsection{~一般概念}
\vsssub

\ws\ 中还提供了一些附加子程序，用作耦合海洋-波浪或海洋-大气系统的一部分。这些子程序用于计算与表面波场有关、并准备传递给外部模式（例如海洋模式）的附加量。设置这些子程序的动机，是让外部模式能够包含波浪对风应力、上层海洋湍流等量的影响。

\paragraph{依赖海况的气海通量}

气海动量通量，或总风应力，是进入表面波和次表层流的动量通量之和。不考虑表面重力波场影响的耦合大气-海洋模式，通常根据风速与风应力之间的经验关系（通过阻力系数 $C_d$）计算总风应力。所提供的 {\it FLD} 子程序允许基于 \ws\ 波数-方向谱计算总风应力，供耦合数值模式使用。

在一阶近似下，总风应力等于进入表面波的动量通量（表面波形状阻力）与直接进入次表层流的动量通量（通过黏性应力）之和。进入波浪的动量通量可表示为波方差谱乘以波增长率（动量吸收率）的积分。计算波形状阻力需要若干假设。首先，波形状阻力在一阶上对高频波（或谱尾）的水平敏感。波谱的这一部分在波浪模式中有很大不确定性，因此在计算风应力时可能需要单独参数化。因此，必须假定一种方式来参数化高频部分；在风速超过 15 m/s 时，该部分不受观测资料约束。其次，需要对波增长率函数作出假设，因为历史上它通常由风速或风应力来参数化。无论哪种情况，增长率函数都需要基于波长以及波向相对于风和/或应力方向的经验系数。第三，波形状阻力会对波边界层（大致为气海界面以上 10 m 内）的湍流剖面和风剖面产生反馈。这种反馈在决定风应力和平均风剖面方面究竟有多重要，目前尚未完全理解。最后，已知破碎波与非破碎波上的增长率不同。然而，目前没有简单方法能在风应力计算模型中显式包含破碎波影响。因此，当前两个 {\it FLD} 子程序均不区分破碎波和非破碎波的增长率。

总气海动量通量可（一阶近似）表示为：
\begin{equation}
\vec{\tau}=\vec{\tau}_{\nu}+\vec{\tau}_{f},
\end{equation}
其中 $\vec{\tau}_{\nu}$ 为黏性应力矢量，$\vec{\tau}_{f}$ 为波形状阻力。在气海界面处，波形状阻力可作为进入所有波浪的动量通量贡献计算：
\begin{equation}\label{formdrag}
\vec{\tau}_{f}=\rho_w\int_{k_{min}}^{k_{max}}\int_{-\pi}^{\pi} \beta_g(k,\theta) \sigma F (k,\theta) d\theta \vec{k} dk,
\end{equation}
其中 $\rho_w$ 为水密度，$k$ 为波数，$\theta$ 为波向，$\sigma$ 为角频率，$\beta_g (k,\theta)$ 为增长率，$F (k,\theta)$ 为波方差谱，$k_{min}$ 和 $k_{max}$ 为产生贡献的波的最小和最大波数。增长率表达式取决于模式所采用的理论，并将在各理论的说明中分别描述。

风速超过 15 m/s 时，谱尾在观测和理论上约束都不足。因此，{\it FLD} 子程序中的谱尾水平以经验方式参数化，使平均阻力系数对应于建模系统中使用的标准整体阻力系数。这样，使用任何显式依赖海况的风应力公式都不会改变风应力的平均值，但应力会根据海况偏离该平均值。这里假定谱尾水平仅是风速的函数，与海况无关。
